\documentclass[11pt,a4paper]{article}

% ============================================================
%  QIMMA -- TEMPLATE "EXERCICES"
%  Charte : logo + dégradé teal/emeraude + accent ambre,
%  fiches-exercices avec onglet de titre, difficulté en étoiles.
% ============================================================

\usepackage[french]{babel}
\usepackage[utf8]{inputenc}
\usepackage[T1]{fontenc}
\usepackage{lmodern}
\usepackage[a4paper,margin=2cm,top=2.3cm,bottom=2.6cm]{geometry}
\usepackage{amsmath,amssymb}
\usepackage{xcolor}
\usepackage{graphicx}
\usepackage{tikz}
\usetikzlibrary{shadows,calc}
\usepackage[most]{tcolorbox}
\usepackage{titlesec}
\usepackage{fancyhdr}
\usepackage{enumitem}
\usepackage{pifont}
\usepackage{hyperref}

% ------------------- CHARTE QIMMA -------------------
\definecolor{qteal}{HTML}{0d9488}
\definecolor{qtealdark}{HTML}{134e4a}
\definecolor{qamber}{HTML}{fbbf24}
\definecolor{qambertext}{HTML}{92400e}
\definecolor{ink}{HTML}{1f2937}
\definecolor{inkgray}{HTML}{6b7280}
\definecolor{tealpale}{HTML}{e6f5f3}
\color{ink}

\graphicspath{{assets/}}
\newcommand{\logofile}{assets/qimma-logo.pdf}

% ------------------- METADONNEES (a completer) -------------------
\newcommand{\matiere}{Matière}                  % TODO
\newcommand{\niveau}{Niveau}                    % TODO
\newcommand{\chapitretitre}{Titre du chapitre}  % TODO

% ------------------- EN-TETE / PIED DE PAGE -------------------
\pagestyle{fancy}
\fancyhf{}
\renewcommand{\headrulewidth}{0pt}
\renewcommand{\footrulewidth}{0.5pt}
\fancyfoot[L]{\raisebox{-0.15cm}{\includegraphics[height=0.35cm]{\logofile}}~\small\sffamily\color{inkgray}Qimma}
\fancyfoot[C]{\small\sffamily\color{inkgray}\matiere\ -- \niveau}
\fancyfoot[R]{\small\sffamily\color{qtealdark}\thepage}

% ------------------- INDICATEUR DE DIFFICULTE -------------------
% Usage : \difficulte{n} avec n = 0..3 étoiles pleines
\newcommand{\difficulte}[1]{%
  \ifcase#1\relax
    {\color{black!25}\ding{72}\ding{72}\ding{72}}%
  \or {\color{qamber}\ding{72}}{\color{black!25}\ding{72}\ding{72}}%
  \or {\color{qamber}\ding{72}\ding{72}}{\color{black!25}\ding{72}}%
  \or {\color{qamber}\ding{72}\ding{72}\ding{72}}%
  \fi
}

% ------------------- BOITE EXERCICE (onglet de titre) -------------------
\newtcolorbox{exercice}[1]{
  enhanced, breakable,
  colback=white, colframe=qteal,
  boxrule=1pt, arc=2mm,
  top=8mm, left=4mm, right=4mm, bottom=4mm,
  drop shadow={black!25, shadow xshift=1mm, shadow yshift=-1mm},
  attach boxed title to top left={xshift=4mm, yshift=-3mm, yshifttext=-1mm},
  boxed title style={colback=qteal, arc=1.5mm, boxrule=0pt, left=2mm, right=2mm},
  coltitle=white, fonttitle=\sffamily\bfseries,
  title={Exercice~#1}
}

\hypersetup{colorlinks=true, linkcolor=qtealdark, urlcolor=qtealdark}

% ============================================================
\begin{document}

% ------------------- BANNIERE DE TITRE (charte Qimma) -------------------
\begin{tcolorbox}[
  enhanced, boxrule=0pt, arc=4mm,
  interior style={left color=qtealdark, right color=qteal},
  top=8mm, bottom=8mm, left=10mm, right=32mm,
  drop shadow={black!35, shadow xshift=1.5mm, shadow yshift=-1.5mm},
  before skip=0pt, after skip=9mm,
  overlay={
    \node[anchor=north west] at ([xshift=6mm,yshift=-6mm]frame.north west)
      {\includegraphics[height=1.25cm]{\logofile}};
    \node[anchor=north east, fill=qamber, text=qambertext, rounded corners=2pt,
          inner sep=4pt, font=\sffamily\bfseries\small] at ([xshift=-6mm,yshift=-6mm]frame.north east) {EXERCICES};
  }
]
\hspace{1.6cm}\centering
{\sffamily\bfseries\color{white}\fontsize{23}{27}\selectfont Série d'exercices}\\[3mm]
{\sffamily\color{white!90}\large \matiere\ \textbullet\ \niveau\ \textbullet\ \chapitretitre}
\end{tcolorbox}

% Légende difficulté
\begin{flushright}
\small\sffamily\color{inkgray}
Difficulté~: {\color{qamber}\ding{72}}~facile~~
{\color{qamber}\ding{72}\ding{72}}~moyen~~
{\color{qamber}\ding{72}\ding{72}\ding{72}}~difficile
\end{flushright}
\vspace{2mm}

% ---------------------------------------------------
\begin{exercice}{1}
\begin{flushright}\difficulte{1}\end{flushright}\vspace{-3mm}
% TODO: énoncé de l'exercice 1

\begin{enumerate}[label=\textbf{\arabic*.}, leftmargin=1cm, itemsep=1mm]
  \item % TODO: question 1
  \item % TODO: question 2
\end{enumerate}
\end{exercice}

\vspace{4mm}

% ---------------------------------------------------
\begin{exercice}{2}
\begin{flushright}\difficulte{2}\end{flushright}\vspace{-3mm}
% TODO: énoncé de l'exercice 2

\begin{enumerate}[label=\textbf{\arabic*.}, leftmargin=1cm, itemsep=1mm]
  \item % TODO: question 1
  \item % TODO: question 2
\end{enumerate}
\end{exercice}

% Dupliquer le bloc \begin{exercice}{numéro} ... \end{exercice} autant que nécessaire.
% \difficulte{0-3} règle le nombre d'étoiles pleines (0 = aucune, 3 = maximum).

% =====================================================
% CORRECTION (optionnelle - à activer si besoin)
% =====================================================
% \newpage
% \begin{tcolorbox}[enhanced, boxrule=0pt, arc=2mm,
%   interior style={left color=qtealdark, right color=qteal},
%   left=6mm, right=6mm, top=4mm, bottom=4mm, fontupper=\color{white}\sffamily,
%   title={\sffamily\bfseries\color{white}Corrigé}, colbacktitle=qtealdark, coltitle=white, boxrule=0pt]
% \subsection*{Exercice 1}
% % TODO: correction exercice 1
% \subsection*{Exercice 2}
% % TODO: correction exercice 2
% \end{tcolorbox}

\end{document}
